\begin{table}[htbp]
  \centering
  \caption{Profile of formal workers by firm-size group, 2025}
  \label{tab:formal-worker-profile-2025}
  \small
  \begin{tabular}{lrrr}
    \toprule
    Characteristic & Formal solo & Formal 2--100 & Formal 101+ \\
    \midrule
    Weighted workers & 494,663 & 2,639,735 & 3,805,111 \\
    Mean real hourly income & 17,217 & 13,602 & 18,548 \\
    Mean age & 44.6 & 38.2 & 37.9 \\
    Women & 45.1\% & 45.1\% & 47.5\% \\
    Higher education & 53.7\% & 52.9\% & 68.0\% \\
    Own-account workers & 83.0\% & 7.2\% & 9.3\% \\
    Domestic workers & 15.8\% & 0.7\% & 0.0\% \\
    Professional and administrative services & 24.2\% & 17.1\% & 11.3\% \\
    Transport and storage & 15.7\% & 4.9\% & 6.0\% \\
    Commerce and repair & 15.3\% & 22.6\% & 11.1\% \\
    Health and social assistance & 5.0\% & 6.3\% & 10.8\% \\
    Households as employers & 15.8\% & 0.7\% & 0.0\% \\
    \bottomrule
  \end{tabular}
  \vspace{0.3em}
  \begin{minipage}{0.95\textwidth}
  \footnotesize
  Notes: Statistics use 2025 formal workers and GEIH expansion weights. The 2--100 column pools all formal workers in firm-size categories from 2--3 through 51--100 workers. Higher education corresponds to workers classified as superior or university educated. Sector rows report the weighted share of each group located in the listed economic activity.
  \end{minipage}
\end{table}
